\section{The closed topological support}\label{sec:support}

The decay assertion in Theorem~\ref{thm:main} holds almost surely under
$Q_{\rm crit}$.  The set of functions that decay at infinity is not closed
in $C_{\rm loc}((0,\infty))$, and in the present problem the distinction is
substantive.  We prove that the closed topological support contains a
nondecaying function.  For an edge configuration $g$, write
$\mathcal A_g$ for the function in \eqref{eq:source-2}.

\begin{proposition}\label{prop:support-constant}
The constant function $s\mapsto 1$ belongs to the closed topological
support of $Q_{\rm crit}$ in $C_{\rm loc}((0,\infty))$.
\end{proposition}

The proof modifies finitely many edge points so that the conditioned law
concentrates near $x_1=\pm1$.  A smooth function distinguishing these
two regions has vanishing form energy, while the resolvent contribution
of the remaining coordinates vanishes.  The finite-shift absolute continuity from
Lemma~\ref{lem:functional-shift} ensures that every function used in the
approximation belongs to the closed support.

\begin{proof}
Let $K=\supp Q_{\rm crit}$ denote the closed topological support.
Since $C_{\rm loc}((0,\infty))$ is separable, $Q_{\rm crit}(K)=1$.
Consider all shifts of finitely many edge points by rational amounts,
restricted to the event that the shifted points remain strictly ordered.
There are countably many such shifts.  By
Lemma~\ref{lem:functional-shift}, outside a single null set, every valid
shift $g'$ of $g$ gives a configuration for which the conditioned measure,
the form identities, and \eqref{eq:source-2} are defined, and for which
$\mathcal A_{g'}$ belongs to $K$.  We fix such a configuration $g$, also
satisfying the spectral bounds of Section~\ref{sec:edge-law}.
Choose and fix $b$ with $g_1+b>0$, and put
\[
 l_a=(g_a+b)^{-1},\qquad D=D(b).
\]
In particular, $l_a$ is strictly decreasing,
$\sum_a l_a=\infty$, and $l_a\le C a^{-2/3}$ for all sufficiently
large $a$.  All constants below may depend on this fixed configuration
and on $b$.

\medskip\noindent
\emph{Choosing the modified variances.}
We leave $g_1$ fixed and modify $g_2,\ldots,g_m$, along a sequence
of integers $m\to\infty$.  Let $J=60$ and define
\[
 A_m=\sum_{i=2}^m l_i+D+l_1-1.
\]
Then $A_m\to\infty$ and $A_m=O(m^{1/3})$.  There are arbitrarily
large $m$ such that
\begin{equation}\label{eq:support-subsequence}
 l_{m+1}\le \frac{A_m}{2m}.
\end{equation}
Indeed, if the reverse strict inequality held eventually, the identity
$A_{m+1}=A_m+l_{m+1}$ would imply
$A_{m+1}>A_m(1+1/(2m))$, and hence $A_m\ge c m^{1/2}$,
contrary to the upper bound.  Restrict henceforth to integers satisfying
\eqref{eq:support-subsequence}, and set $M=A_m/(m-1)$.  For all
sufficiently large such integers,
\begin{equation}\label{eq:support-scale}
 \frac{c}{m}\le M\le C m^{-2/3}.
\end{equation}

Choose strictly decreasing positive numbers on indices $2,\ldots,m$
whose sum is $A_m$, with the first $J-1$ numbers close to $2M$ and the
others close to $M$.  This is possible by adjusting the latter group
by a relative amount $O(J/m)$ and then making small distinct
perturbations of zero total sum.  For large $m$ these numbers are less
than $l_1$ and greater than $l_{m+1}$, by
\eqref{eq:support-subsequence}.  They therefore define an ordered
modification of the edge points through $g_i'=1/l_i'-b$.
Approximate each of the finitely many shifts $g_i'-g_i$ by a rational
number, sufficiently closely to retain strict order and all the strict
comparisons just imposed.  We continue to write $g^{(m)}$ and $l_a^{(m)}$
for the resulting configuration and variances.  The approximation can
also be chosen so that the following bounds hold, with constants
independent of $m$:
\begin{align}
 &l_1^{(m)}=l_1,\qquad l_a^{(m)}=l_a\quad(a>m),
       \qquad l_a^{(m)}\asymp M\quad(2\le a\le m),
       \label{eq:support-variances}\\
 &\frac{l_a^{(m)}}{l_j^{(m)}}\le q<1
       \quad(a>J,\ 1\le j\le J),
       \label{eq:support-separation}\\
 &V_m:=D^{(m)}(b)+l_1\longrightarrow1,
       \qquad
       \sum_{a>1}(l_a^{(m)})^2\le C m^{-1/3}.
       \label{eq:support-square-sum}
\end{align}
Here $D^{(m)}$ is the value of \eqref{eq:source-1} for $g^{(m)}$.
To check the first conclusion in \eqref{eq:support-square-sum}, the
finite-change identity from that equation gives
\[
 D^{(m)}(b)-D(b)
   =-\sum_{i=2}^m(l_i^{(m)}-l_i).
\]
Before rational approximation our choice of sum makes $V_m=1$ exactly;
afterwards we can require $|V_m-1|\le m^{-1}$.
The square-sum bound follows from $mM^2\le C m^{-1/3}$ and
$\sum_{a>m}l_a^2\le C m^{-1/3}$.  By the choice of the original
configuration, each $\mathcal A_{g^{(m)}}$ belongs to $K$.

\medskip\noindent
\emph{Concentration after conditioning.}
Let $\gamma_m$ be the product law of centered normals on indices $a>1$
with variances $l_a^{(m)}$, and let
\[
 S_m=\sum_{a>1}\bigl(x_a^2-l_a^{(m)}\bigr).
\]
This series converges in $L^2(\gamma_m)$, and
\eqref{eq:support-square-sum} implies $S_m\to0$ in probability under
$\gamma_m$.  Write $\Pi_m=\Pi_{g^{(m)}}$.  Eliminating $x_1$ in the
conditioning prescription shows that the $a>1$ marginal of $\Pi_m$
has density
\begin{equation}\label{eq:support-tail-density}
 \begin{aligned}
 r_m&=Z_m^{-1}h_m(S_m),\qquad Z_m=\E_{\gamma_m}h_m(S_m),\\
 h_m(v)&=\1_{\{v<V_m\}}(V_m-v)^{-1/2}
             \exp\left(-\frac{V_m-v}{2l_1}\right),
 \end{aligned}
\end{equation}
relative to $\gamma_m$.  Under the full law, $x_1^2=V_m-S_m$,
and the sign of $x_1$ is independent and symmetric.

We give bounds near the singularity in
\eqref{eq:support-tail-density}, because concentration under the
product law alone would not suffice.  Put $t_m=m^{1/10}$.  For either
sign, the centered Gaussian square-sum formula and
\eqref{eq:support-square-sum} give
\begin{equation}\label{eq:support-mgf}
 \log\E_{\gamma_m}e^{\pm t_m S_m}
 \le C t_m^2\sum_{a>1}(l_a^{(m)})^2\le C,
\end{equation}
since $t_m\max_{a>1}l_a^{(m)}\to0$.
Under either exponential tilt the coordinates remain independent
normals, with variances
$l_a^{(m)}/(1\mp2t_m l_a^{(m)})$, comparable to the original
variances.  The sum of the squares of any fixed four coordinates among
$2,\ldots,J$ has density bounded by $C/M$: its characteristic function
has absolute value at most $C(1+M^2u^2)^{-1}$, whose integral is
$O(M^{-1})$.  Convolution with the remaining coordinates does not
increase this bound.  The same assertion holds for the infinite sum
by its $L^2$ convergence.  Undoing the tilt and using
\eqref{eq:support-scale}--\eqref{eq:support-mgf}, we obtain a density
$p_m$ of $S_m$ satisfying
\begin{equation}\label{eq:support-density-bound}
 p_m(v)\le C m e^{-t_m|v|},\qquad v\in\R.
\end{equation}

The factors $h_m(v)$ are uniformly bounded on the complement of a fixed
neighborhood of $v=V_m$.  Near that point, which stays bounded away
from zero, \eqref{eq:support-density-bound} supplies an exponentially
small factor, while the singularity $(V_m-v)^{-p/2}$ is integrable
for every $p<2$.  Consequently, concentration of $S_m$ at zero gives
\begin{equation}\label{eq:support-density-norm}
 Z_m\longrightarrow e^{-1/(2l_1)},\qquad
 \sup_m\|r_m\|_{L^{3/2}(\gamma_m)}<\infty.
\end{equation}
The same estimate, integrated over $|S_m|>\varepsilon$, shows that
for each fixed $\varepsilon>0$ there are positive constants
$C_\varepsilon,c_\varepsilon$ such that
\begin{equation}\label{eq:support-concentration}
 \Pi_m\{|x_1^2-1|>\varepsilon\}
       \le C_\varepsilon e^{-c_\varepsilon m^{1/10}}
\end{equation}
for all sufficiently large $m$.  Integrating the negative tail of
$S_m$ also gives convergence of $x_1^2$ to $1$ in every fixed moment.
Finally, H\"older's inequality applied to
\eqref{eq:support-density-norm} yields, for each fixed $p<\infty$,
\begin{equation}\label{eq:support-coordinate-moments}
 \sup_{m}\sup_{a>1}
 \Pi_m\left[\left|x_a/\sqrt{l_a^{(m)}}\right|^p\right]<\infty.
\end{equation}

\medskip\noindent
\emph{The first coordinate and the trace tail.}
Write $H_m=H_{\Pi_m}$ and
$\mathcal E_m(F)=\int|\nabla F|^2\,d\Pi_m$.
Choose an odd smooth compactly supported function $\phi$ on $\R$
which equals $1$ near $1$ and $-1$ near $-1$.  The preceding
concentration and moment estimates imply
\[
 \|x_1-\phi(x_1)\|_{L^2(\Pi_m)}\longrightarrow0,\qquad
 \|\phi(x_1)\|_{L^2(\Pi_m)}^2\longrightarrow1,\qquad
 \mathcal E_m(\phi(x_1))\longrightarrow0.
\]
For $0\le s\le S$, the spectral theorem gives
\[
 0\le \|\phi(x_1)\|_2^2
       -\langle\phi(x_1),e^{-s c_*H_m}\phi(x_1)\rangle
       \le S c_*\mathcal E_m(\phi(x_1)).
\]
Semigroup contraction therefore shows that the first coordinate's
contribution to $\mathcal A_{g^{(m)}}(s)$ tends to $1$, uniformly
for $s$ in bounded intervals.  Coordinates $2,\ldots,J$ contribute
at most $\sum_{a=2}^J\Pi_m[x_a^2]=O(M)$, which tends to zero.
It remains to control the infinitely many coordinates beyond $J$.
We use the angular integration-by-parts estimate of
Section~\ref{sec:trace}.  Here the angular identity
\eqref{eq:source-23} is exact, so there is no likelihood-score error.
The remaining task is to make the coefficient and inverse-moment
bounds uniform in $m$.

Let $R_0=\sum_{j=1}^Jx_j^2$.  We first establish the uniform inverse
moments needed for angular integration by parts.  By
\eqref{eq:support-density-norm}, H\"older's inequality, and
$R_0\ge\sum_{j=2}^Jx_j^2$,
\[
 \Pi_m[R_0^{-8}]
 \le C\left(\E_{\gamma_m}
       \left[\left(\sum_{j=2}^Jx_j^2\right)^{-24}\right]\right)^{1/3}
 \le C M^{-8}\le C m^8.
\]
The Gaussian inverse moment is finite because $J-1=59>48$.
On $\{x_1^2\ge1/2\}$ the quantity $R_0^{-4}$ is bounded by $16$;
on the complement, Cauchy--Schwarz and
\eqref{eq:support-concentration} absorb the preceding polynomial
bound.  Thus
\begin{equation}\label{eq:support-inverse-moments}
 \sup_m\Pi_m[R_0^{-4}]<\infty.
\end{equation}

For $a>J$ put
\[
 \delta_{aj}=(l_a^{(m)})^{-1}-(l_j^{(m)})^{-1},\qquad
 \mathcal V_a=\sum_{j=1}^J\frac{x_j}{R_0\delta_{aj}}
                   (x_j\partial_a-x_a\partial_j).
\]
By \eqref{eq:support-separation},
$\delta_{aj}\ge(1-q)/l_a^{(m)}$.  The calculation in
\eqref{eq:source-42}, now in the limiting coordinates, gives
\begin{align*}
 |\mathcal V_a|
   &\le C l_a^{(m)}(1+|x_a|/\sqrt{R_0}),\\
 |\operatorname{div}\mathcal V_a|
   &\le C l_a^{(m)}|x_a|/R_0.
\end{align*}
Together with \eqref{eq:support-coordinate-moments} and
\eqref{eq:support-inverse-moments}, these imply
\begin{equation}\label{eq:support-field-norms}
 \|\mathcal V_a\|_{L^2(\Pi_m)}
 +\|\operatorname{div}\mathcal V_a\|_{L^2(\Pi_m)}
 \le C l_a^{(m)}.
\end{equation}

Each rotation $x_j\partial_a-x_a\partial_j$ preserves the square-sum
constraint.  Its derivative of the Gaussian log density is
$-\delta_{aj}x_ax_j$, so the corresponding derivative along
$\mathcal V_a$ is exactly $-x_a$.  Applying
\eqref{eq:source-23} with the coefficients defining $\mathcal V_a$
therefore gives
\begin{equation}\label{eq:support-ibp}
 \Pi_m[x_af]=\Pi_m[\mathcal V_a f]
                   +\Pi_m[f\operatorname{div}\mathcal V_a]
\end{equation}
for smooth cylinder tests.  To justify the singular coefficients,
multiply them first by $\chi(R_0/\varepsilon)$, where $\chi$ is smooth,
zero below $1$, and one above $2$.  Direct calculation gives
\[
 \mathcal V_aR_0
  =-2x_a\sum_{j=1}^J\frac{x_j^2}{R_0\delta_{aj}},\qquad
 |\mathcal V_a\chi(R_0/\varepsilon)|
  \le C l_a^{(m)}\frac{|x_a|}{R_0}
           \1_{\{R_0\le2\varepsilon\}}.
\]
The last expression tends to zero in $L^2(\Pi_m)$ by
\eqref{eq:support-coordinate-moments} and
\eqref{eq:support-inverse-moments}.  At infinity we use a radial cutoff
in a finite coordinate set containing $1,\ldots,J,a$ and all test
coordinates.  Every rotation in $\mathcal V_a$ preserves this radius,
so differentiating that cutoff produces no error.  This proves
\eqref{eq:support-ibp} on the cylinder core, which then extends the
identity to the form domain.

Fix $z>0$ and a finite set $I\subset\{J+1,J+2,\ldots\}$.  Applying
\eqref{eq:support-ibp} to a vector of tests $(f_a)_{a\in I}$ and
using Cauchy--Schwarz and \eqref{eq:support-field-norms} yields
\[
 \left|\sum_{a\in I}\Pi_m[x_af_a]\right|
 \le C_z\left(\sum_{a>J}(l_a^{(m)})^2\right)^{1/2}
       \left(\sum_{a\in I}
          \{z\|f_a\|_2^2+c_*\mathcal E_m(f_a)\}\right)^{1/2}.
\]
By the resolvent variational formula, followed by monotone convergence
over finite $I$,
\begin{equation}\label{eq:support-resolvent-tail}
 \sum_{a>J}\langle x_a,(z+c_*H_m)^{-1}x_a\rangle_{\Pi_m}
       \le C_z\sum_{a>J}(l_a^{(m)})^2\longrightarrow0.
\end{equation}
For every $s_0>0$ the scalar inequality
$e^{-s\lambda}\le C_{s_0,z}(z+\lambda)^{-1}$ holds for
$s\ge s_0$ and $\lambda\ge0$.  Hence
\eqref{eq:support-resolvent-tail} also makes the semigroup trace
beyond $J$ tend to zero uniformly on compact subsets of $(0,\infty)$.

We have proved $\mathcal A_{g^{(m)}}\to1$ in
$C_{\rm loc}((0,\infty))$.  Each approximating function belongs to
$K$, so closedness of $K$ proves the proposition.
\end{proof}
